\begin{example} % 直线，三点共线
    % 斜率例题
（2026朝阳期末）已知椭圆 $E:\frac{x^2}{a^2}+\frac{y^2}{b^2}=1\ (a>b>0)$ 的离心率为 $\frac{\sqrt{2}}{2}$，椭圆 $E$ 上的点到两焦点的距离之和为 $2\sqrt{6}$.\\
（1）求椭圆 $E$ 的方程；\\
（2）设点 $M(m,0),N(n,0)\ (m\neq n)$，过点 $M$ 的直线 $l$ 与椭圆 $E$ 交于不同的两点 $A,B$，点 $A$ 关于 $x$ 轴的对称点记为 $P$（$P$ 与 $B$ 不重合）. 若 $mn=6$，试判断 $P,B,N$ 三点是否共线，并说明理由.
\end{example}

\begin{example} % 直线，直角三角形特殊角
    % 特殊三角形例题
（2026西城期末）已知椭圆 $E:\frac{x^2}{2}+y^2=1$ 的上顶点为 $A$，动点 $B$ 为 $x$ 轴正半轴上的一点，过点 $A$ 作 $AB$ 的垂线交椭圆 $E$ 于另一点 $C$.\\
（1）当直线 $AC$ 经过椭圆的左焦点时，求点 $B$ 的坐标；\\
（2）问是否存在点 $B$，分别使得 $\angle ACB$ 等于 $30^\circ,45^\circ$ 和 $60^\circ$？说明理由.
\end{example}

\begin{example} % 直线或者动点，距离公式，等腰直角三角形
    %特殊三角形例题，
    % 距离例题
（2026海淀期末）已知椭圆 $E:\frac{x^2}{a^2}+\frac{y^2}{b^2}=1\ (a>b>0)$ 的左顶点为 $A$，请从下面条件中任选两个作为已知.\\
条件①：椭圆 $E$ 的离心率为 $\frac{1}{2}$；\\
条件②：椭圆 $E$ 经过点 $(0,\sqrt{3})$；\\
条件③：点 $A$ 到椭圆 $E$ 的左焦点的距离为 $1$.\\
（1）求椭圆 $E$ 的方程；\\
（2）过点 $A$ 的直线 $l$ 与椭圆 $E$ 的另一个交点为 $P$，$l$ 与 $y$ 轴交于点 $M$. 过点 $P$ 作 $l$ 的垂线交直线 $x=2$ 于点 $N$. 若 $\triangle PMN$ 是等腰直角三角形，求 $l$ 的斜率.
\end{example}

\begin{example} % 动点，韦达定理求坐标
（2026东城期末）已知椭圆 $C:\frac{x^2}{a^2}+\frac{y^2}{b^2}=1\ (a>b>0)$ 的左、右焦点分别为 $F_1,F_2$，离心率为 $\frac{\sqrt{5}}{5}$. 设 $P(m,n)$ 是椭圆 $C$ 上一点，且在第一象限，$\lvert PF_1\rvert+\lvert PF_2\rvert=2\sqrt{5}$.\\
（1）求椭圆 $C$ 的方程；\\
（2）已知 $PF_2$ 不与 $x$ 轴垂直，延长 $PF_1,PF_2$ 分别与椭圆 $C$ 交于点 $M,N$. 若点 $M$ 到直线 $x=-3$ 的距离与点 $N$ 到直线 $x=3$ 的距离之和为 $\frac{24}{5}$，求 $m$ 的值.
\end{example}

\begin{example} % 动点，直线相交求坐标，距离公式
% 入题思路例题1.4
% 距离例题
（2026丰台期末）已知 $F_1,F_2$ 是椭圆 $E:\frac{x^2}{a^2}+\frac{y^2}{b^2}=1\ (a>b>0)$ 的左、右焦点，且 $\lvert F_1F_2\rvert=2$，点 $Q\left(1,\frac{3}{2}\right)$ 在椭圆 $E$ 上.\\
（1）求椭圆 $E$ 的标准方程；\\
（2）设 $P(x_0,y_0)$ 为第一象限内 $E$ 上的动点（不与 $Q$ 重合），直线 $PF_1,PF_2$ 分别交直线 $l:4y_0x-3x_0y=0$ 于点 $M,N$. 求 $\lvert MF_1\rvert+\lvert NF_2\rvert$ 的值.
\end{example}

\begin{example} % 直线、正方形翻译
    % 平行四边形例题
（2025西城期末）已知椭圆 $E:\frac{x^2}{a^2}+\frac{y^2}{b^2}=1\ (a>b>0)$ 的左、右顶点分别为 $A_1,A_2$，离心率为 $\frac{\sqrt{2}}{2}$，点 $T(0,1)$，$\triangle TA_1A_2$ 的面积为 $2$.\\
（1）求椭圆 $E$ 的方程；\\
（2）过点 $T$ 且斜率为 $k$ 的直线交椭圆 $E$ 于点 $C,D$，线段 $CD$ 的垂直平分线交 $y$ 轴于点 $Q$，点 $Q$ 关于直线 $CD$ 的对称点为 $P$. 若四边形 $PCQD$ 为正方形，求 $k$ 的值.
\end{example}

\begin{example} % 动点，韦达定理求坐标，平行关系
    % 练习题
（2025海淀期末）已知椭圆 $C:\frac{x^2}{a^2}+\frac{y^2}{b^2}=1\ (a>b>0)$ 的左顶点为 $A(-2,0)$，离心率 $e=\frac{\sqrt{3}}{2}$.\\
（1）求 $C$ 的标准方程；\\
（2）设点 $P$ 为 $C$ 上异于顶点的一点，点 $P$ 关于 $y$ 轴的对称点为 $Q$，过 $A$ 作 $OP$ 的平行线 $l$，$l$ 与 $C$ 的另一个交点为 $M$. 当 $M$ 与 $Q$ 不重合时，求证：$MQ\parallel AP$.
\end{example}

\begin{example} % 动点，韦达定理求坐标，距离
    % 练习题
（2025东城期末）已知椭圆 $E:\frac{x^2}{a^2}+\frac{y^2}{b^2}=1\ (a>b>0)$ 的一个顶点为 $(2,0)$，且焦距为 $2$，$A$ 为第一象限内 $E$ 上的动点. 过点 $A$ 作斜率为 $\frac{1}{2},-\frac{1}{2}$ 的直线分别与 $E$ 交于点 $M,N$（均异于点 $A$），直线 $MN$ 与 $x$ 轴交于点 $H$，$Q$ 为线段 $MN$ 的中点，直线 $AQ$ 与 $x$ 轴交于点 $P$.\\
（1）求 $E$ 的方程；\\
（2）当 $\lvert HP\rvert=3$ 时，求点 $A$ 的横坐标.
\end{example}

\begin{example} % 入题思路的选择，直曲联立求坐标
% 入题思路例题1.11
（2025丰台期末）已知椭圆 $C:\frac{x^2}{a^2}+\frac{y^2}{b^2}=1\ (a>b>0)$ 的上顶点为 $A(0,1)$，离心率为 $\frac{\sqrt{3}}{2}$.\\
（1）求椭圆 $C$ 的方程；\\
（2）设 $B$ 为椭圆 $C$ 的下顶点，动点 $M$ 到坐标原点 $O$ 的距离等于 $1$（$M$ 与 $A,B$ 不重合），直线 $AM$ 与椭圆 $C$ 的另一个交点为 $N$. 记直线 $BM,BN$ 的斜率分别为 $k_1,k_2$. 问：是否存在常数 $\lambda$，使得 $k_1+\lambda k_2=0$ 恒成立？若存在，求出 $\lambda$ 的值；若不存在，说明理由.
\end{example}

\begin{example} % 动点（题目误导），韦达定理求坐标，定值定点（特殊情况先得答案）
% 入题思路例题1.13
（2025朝阳期末）已知椭圆 $E:\frac{x^2}{a^2}+\frac{y^2}{b^2}=1\ (a>b>0)$ 的离心率为 $\frac{\sqrt{6}}{3}$，右顶点为 $(2,0)$.\\
（1）求椭圆 $E$ 的方程；\\
（2）过原点 $O$ 且与 $y$ 轴不重合的直线 $l$ 与椭圆 $E$ 交于 $M,N$ 两点. 已知点 $P(0,2)$，直线 $PM,PN$ 与椭圆 $E$ 的另一个交点分别为 $A,B$. 证明：直线 $AB$ 过定点.
\end{example}

\begin{example} % 动点，韦达定理求坐标，三点共线
    % 练习题
（2024东城期末）已知椭圆 $C:\frac{x^2}{a^2}+\frac{y^2}{b^2}=1\ (a>b>0)$ 的右焦点为 $F$，左、右顶点分别为 $A,B$，$\lvert AF\rvert=2+\sqrt{3}$，$\lvert BF\rvert=2-\sqrt{3}$.\\
（1）求椭圆 $C$ 的方程；\\
（2）设 $O$ 是坐标原点，$M,N$ 是椭圆 $C$ 上不同的两点，且关于 $x$ 轴对称，$E,G$ 分别为线段 $OM,MB$ 的中点，直线 $AE$ 与椭圆 $C$ 交于另一点 $D$. 证明：$D,G,N$ 三点共线.
\end{example}

\begin{example} % 直线，基础翻译
% 入题思路例题1.8
（2024朝阳期末）已知椭圆 $E:\frac{x^2}{a^2}+\frac{y^2}{b^2}=1\ (a>b>0)$ 的左顶点为 $A$，上顶点为 $B$，原点 $O$ 到直线 $AB$ 的距离为 $\frac{2\sqrt{5}}{5}$，$\triangle AOB$ 的面积为 $1$.\\
（1）求椭圆 $E$ 的方程；\\
（2）过点 $P(-2,1)$ 的直线 $l$ 与椭圆 $E$ 交于不同的两点 $C,D$，过点 $C$ 作 $x$ 轴的垂线，分别与直线 $AD,AB$ 交于点 $M,N$. 判断点 $N$ 是否为线段 $CM$ 的中点，说明理由.
\end{example}

\begin{example} % 直线，与椭圆相关的中点的处理
% 条件翻译中点处理
（2024西城期末）已知椭圆 $E:\frac{x^2}{a^2}+\frac{y^2}{b^2}=1\ (a>b>0)$ 的离心率为 $\frac{\sqrt{3}}{2}$，且经过点 $C(2,1)$.\\
（1）求 $E$ 的方程；\\
（2）过点 $N(0,1)$ 的直线交 $E$ 于点 $A,B$（点 $A,B$ 与点 $C$ 不重合）. 设 $AB$ 的中点为 $M$，连接 $CM$ 并延长交 $E$ 于点 $D$. 若 $M$ 恰为 $CD$ 的中点，求直线 $AB$ 的方程.
\end{example}

\begin{example} % 直线，直线相交求坐标，求最值
% 入题思路例题1.7
（2024海淀期末）已知椭圆 $E:\frac{x^2}{a^2}+\frac{y^2}{b^2}=1\ (a>b>0)$ 过点 $A(3,0)$，焦距为 $2\sqrt{5}$.\\
（1）求椭圆 $E$ 的方程，并求其短轴长；\\
（2）过点 $P(1,0)$ 且不与 $x$ 轴重合的直线 $l$ 交椭圆 $E$ 于两点 $C,D$，连接 $CO$ 并延长交椭圆 $E$ 于点 $M$，直线 $AM$ 与 $l$ 交于点 $N$，$Q$ 为 $OD$ 的中点，其中 $O$ 为原点. 设直线 $NQ$ 的斜率为 $k$，求 $k$ 的最大值.
\end{example}

\begin{example} % 直线动点，椭圆切线
    % 椭圆切线例题
（2024丰台期末）已知椭圆 $E:\frac{x^2}{4}+\frac{y^2}{3}=1$.\\
（1）求椭圆 $E$ 的离心率和焦点坐标；\\
（2）设直线 $l_1:y=kx+m$ 与椭圆 $E$ 相切于第一象限内的点 $P$，不过原点 $O$ 且平行于 $l_1$ 的直线 $l_2$ 与椭圆 $E$ 交于不同的两点 $A,B$，点 $A$ 关于原点 $O$ 的对称点为 $C$. 记直线 $OP$ 的斜率为 $k_1$，直线 $BC$ 的斜率为 $k_2$，求 $\frac{k_1}{k_2}$ 的值.
\end{example}

\begin{example} % 直线，基础题型
（2023朝阳期末）已知椭圆 $C:\frac{x^2}{a^2}+\frac{y^2}{b^2}=1\ (a>b>0)$ 的右顶点为 $A(2,0)$，$P$ 为椭圆 $C$ 上的动点，且点 $P$ 不在 $x$ 轴上，$O$ 是坐标原点，$\triangle AOP$ 面积的最大值为 $1$.\\
（1）求椭圆 $C$ 的方程及离心率；\\
（2）过点 $H(-1,0)$ 的直线 $PH$ 与椭圆 $C$ 交于另一点 $Q$，直线 $AP,AQ$ 分别与 $y$ 轴相交于点 $E,F$. 当 $\lvert EF\rvert=2$ 时，求直线 $PH$ 的方程.
\end{example}

\begin{example} % 动点，求最值
（2023东城期末）已知椭圆 $C:\frac{x^2}{a^2}+\frac{y^2}{b^2}=1\ (a>b>0)$ 的离心率为 $\frac{\sqrt{3}}{2}$，长轴长与短轴长的和为 $6$，$F_1,F_2$ 分别为椭圆 $C$ 的左、右焦点.\\
（1）求椭圆 $C$ 的方程；\\
（2）设 $P$ 为椭圆 $C$ 上一点，$M(1,0)$. 若 $\lvert PF_1\rvert,\lambda\lvert PM\rvert,\lvert PF_2\rvert$ 成等差数列，求实数 $\lambda$ 的范围.
\end{example}

\begin{example} % 直线，基础题型
（2023海淀期末）已知椭圆 $E:\frac{x^2}{a^2}+\frac{y^2}{b^2}=1$ 过点 $P(-2,1)$ 和 $Q(2\sqrt{2},0)$.\\
（1）求椭圆 $E$ 的方程；\\
（2）过点 $G(0,2)$ 作直线 $l$ 交椭圆 $E$ 于不同的两点 $A,B$，直线 $PA$ 交 $y$ 轴于点 $M$，直线 $PB$ 交 $y$ 轴于点 $N$. 若 $\lvert GM\rvert\cdot\lvert GN\rvert=2$，求直线 $l$ 的方程.
\end{example}

\begin{example} % 动点（在直线上），基础题型
% 练习题目
（2023丰台期末）已知椭圆 $E:\frac{x^2}{a^2}+\frac{y^2}{b^2}=1\ (a>b>0)$ 过点 $A(-2,0)$，离心率为 $\frac{\sqrt{2}}{2}$.\\
（1）求椭圆 $E$ 的方程；\\
（2）设点 $P(2,m)\ (m>0)$，直线 $PA$ 与椭圆 $E$ 的另一个交点为 $C$，$O$ 为坐标原点，$B$ 为椭圆 $E$ 的右顶点. 记直线 $OP$ 的斜率为 $k_1$，直线 $BC$ 的斜率为 $k_2$，求证：$k_1k_2$ 为定值.
\end{example}

\begin{example} % 直线，与椭圆相关的中点处理
% 条件翻译中点处理
（2023西城期末）如图，已知椭圆 $E:\frac{y^2}{a^2}+\frac{x^2}{b^2}=1\ (a>b>0)$ 的一个焦点为 $F_1(0,1)$，离心率为 $\frac{\sqrt{2}}{2}$.\\
（1）求椭圆 $E$ 的方程；\\
（2）过点 $F_1$ 作斜率为 $k$ 的直线交椭圆 $E$ 于两点 $A,B$，$AB$ 的中点为 $M$. 设 $O$ 为原点，射线 $OM$ 交椭圆 $E$ 于点 $C$. 当 $\triangle ABC$ 与 $\triangle ABO$ 的面积相等时，求 $k$ 的值.
\end{example}

\begin{example} % 直线，基础题型
（2022丰台期末）已知椭圆 $C:\frac{x^2}{a^2}+\frac{y^2}{b^2}=1\ (a>b>0)$ 过点 $(\sqrt{2},1)$，离心率为 $\frac{\sqrt{2}}{2}$.\\
（1）求椭圆 $C$ 的方程；\\
（2）设椭圆 $C$ 的右顶点为 $A$，过点 $D(4,0)$ 的直线 $l$ 与椭圆 $C$ 交于不同的两点 $M,N$（均异于点 $A$），直线 $AM,AN$ 分别与直线 $x=4$ 交于点 $P,Q$. 求证：$\lvert DP\rvert\cdot\lvert DQ\rvert$ 为定值.
\end{example}

\begin{example} % 直线，基础题型
（2022朝阳期末）已知曲线 $W:\frac{x^2}{3-m}+\frac{y^2}{m}=1\ (m\in\mathbb{R},\ m\neq0,\ m\neq3)$.\\
（1）若曲线 $W$ 是焦点在 $x$ 轴上的椭圆，求 $m$ 的取值范围；\\
（2）当 $m=1$ 时，过点 $E(1,0)$ 作斜率为 $k\ (k\neq0)$ 的直线 $l$ 交曲线 $W$ 于点 $A,B$（$A,B$ 异于顶点），交直线 $x=2$ 于 $P$. 过点 $P$ 作 $y$ 轴的垂线，垂足为 $Q$，直线 $AQ$ 交 $x$ 轴于 $C$，直线 $BQ$ 交 $x$ 轴于 $D$，求线段 $CD$ 中点 $M$ 的坐标.
\end{example}

\begin{example} % 动点，基础题型
（2022东城期末）已知椭圆 $C:\frac{x^2}{a^2}+\frac{y^2}{b^2}=1$ 过点 $A(-\sqrt{3},0)$，其右焦点为 $F(1,0)$.\\
（1）求椭圆 $C$ 的方程；\\
（2）设 $P$ 为椭圆 $C$ 上一动点（不在 $x$ 轴上），$M$ 为 $AP$ 中点，过原点 $O$ 作 $AP$ 的平行线，与直线 $x=3$ 交于点 $Q$. 问：直线 $OM$ 与 $FQ$ 斜率的乘积是否为定值？若为定值，求出该值；若不为定值，请说明理由.
\end{example}

\begin{example} % 直线，基础题型
（2022海淀期末）已知点 $A(0,-1)$ 在椭圆 $C:\frac{x^2}{3}+\frac{y^2}{b^2}=1$ 上.\\
（1）求椭圆 $C$ 的方程和离心率；\\
（2）设直线 $l:y=k(x-1)\ (k\neq1)$ 与椭圆 $C$ 交于不同两点 $E,F$，直线 $AE,AF$ 分别交直线 $x=3$ 于点 $M,N$. 当 $\triangle AMN$ 的面积为 $3\sqrt{3}$ 时，求 $k$ 的值.
\end{example}
\begin{example} % 直线，铅锤法求面积，非对称韦达
    % 练习题
（2022西城期末）已知椭圆$M: \frac{x^2}{a^2} + \frac{y^2}{b^2} = 1(a > b > 0)$的焦点为$F(2, 0)$，长轴长与短轴长的比值为$\sqrt{2}$.
（1）求椭圆$M$的方程；\\
（2）过点$F$的直线交椭圆$M$于$A, B$两点，$BC \prep x$轴于点$C$，$AD \prep x$轴于点$D$，直线$BD$交直线$x = 4$于点$E$，求$\triangle ECD$与$\triangle EAB$的面积之比.
\end{example}


\begin{example} % 动点，三角形面积（这题不用选）
（2021海淀期末）已知椭圆 $W:\frac{x^2}{a^2}+\frac{y^2}{b^2}=1\ (a>b>0)$ 的离心率为 $\frac{\sqrt{3}}{2}$，且经过点 $C(2,\sqrt{3})$.\\
（1）求椭圆 $W$ 的方程及其长轴长；\\
（2）$A,B$ 分别为椭圆 $W$ 的左、右顶点，点 $D$ 在椭圆 $W$ 上，且位于 $x$ 轴下方，直线 $CD$ 交 $x$ 轴于点 $Q$. 若 $\triangle ACQ$ 的面积比 $\triangle BDQ$ 的面积大 $2\sqrt{3}$，求点 $D$ 的坐标.
\end{example}

\begin{example} % 直线，弦长公式，取值范围
    % 练习题
（2021朝阳期末）已知椭圆 $C:\frac{x^2}{a^2}+\frac{y^2}{b^2}=1\ (a>b>0)$ 过点 $\left(1,\frac{\sqrt{3}}{2}\right)$，且 $C$ 的离心率为 $\frac{\sqrt{3}}{2}$.\\
（1）求椭圆 $C$ 的方程；\\
（2）过点 $P(1,0)$ 的直线 $l$ 交椭圆 $C$ 于 $A,B$ 两点，求 $\lvert PA\rvert\cdot\lvert PB\rvert$ 的取值范围.
\end{example}

\begin{example} % 直线引入，等腰直角三角形的翻译，与直线无关的量的处理
    % 练习题
（2021西城期末）已知椭圆 $C:\frac{x^2}{4}+\frac{y^2}{2}=1$.\\
（1）求椭圆 $C$ 的离心率和长轴长；\\
（2）已知直线 $y=kx+2$ 与椭圆 $C$ 有两个不同的交点 $A,B$，$P$ 为 $x$ 轴上一点. 是否存在实数 $k$，使得 $\triangle PAB$ 是以点 $P$ 为直角顶点的等腰直角三角形？若存在，求出 $k$ 的值及点 $P$ 的坐标；若不存在，说明理由.
\end{example}

\begin{example} % 椭圆切线，距离公式
    % 椭圆切线例题
（2021东城期末）已知椭圆 $C:\frac{x^2}{a^2}+\frac{y^2}{b^2}=1\ (a>b>0)$ 过点 $A(-2,0),B(2,0)$，且离心率为 $\frac{1}{2}$.\\
（1）求椭圆 $C$ 的方程；\\
（2）设直线 $l$ 与椭圆 $C$ 有且仅有一个公共点 $E$，且与 $x$ 轴交于点 $G$（$E,G$ 不重合），$ET\perp x$ 轴，垂足为 $T$. 求证：$\frac{\lvert TA\rvert}{\lvert TB\rvert}=\frac{\lvert GA\rvert}{\lvert GB\rvert}$.
\end{example}

\begin{example} % 直线，定值，特殊情况先得答案，非对称韦达
    % 定值非对称例题
（2021丰台期末）已知椭圆 $W:\frac{x^2}{a^2}+\frac{y^2}{b^2}=1\ (a>b>0)$ 过 $A(0,2),B(-3,-1)$ 两点.\\
（1）求椭圆 $W$ 的方程；\\
（2）直线 $AB$ 与 $x$ 轴交于点 $M(m,0)$，过点 $M$ 作不垂直于坐标轴且与 $AB$ 不重合的直线 $l$，$l$ 与椭圆 $W$ 交于 $C,D$ 两点，直线 $AC,BD$ 分别交直线 $x=m$ 于 $P,Q$ 两点，求证：$\frac{\lvert PM\rvert}{\lvert MQ\rvert}$ 为定值.
\end{example}
